% a07_k01steps_c: last part of app:a07_k01steps (input from sec/a07_k01steps.tex): the regrouping into full classes of the cut components, their coverage by Balaban's renormalisation bound, and the theorem.

\begin{lemma}[the expansion]\label{lem:a07:M}
Put $G_C\defeq G_k(\OmC,\Bt)$ ($=G_k(\OmC,B^{(k+1),\eta})$ by \cref{lem:a07:C51}), $\dG\defeq G_A-G_C$ on $\OmC\times\OmC$ with $G_C$ extended by $0$, $E_C\defeq a_kG_Cf$ and $\dE\defeq E_A-E_C=a_k\dG f_{\mathrm{in}}+a_kG_Af_{\mathrm{out}}$ with $f_{\mathrm{in}}\defeq\mathbf 1_{\OmC}f$, $f_{\mathrm{out}}\defeq f-f_{\mathrm{in}}$. For a class $\cC$, or a product of classes contracted by fixed fields,
\[
\cC[G_A;E_A]-\cC[G_C;E_C]\ =\ \sum_{(n,m)\ne(0,0)}\cC^{(n,m)},
\]
where $\cC^{(n,m)}$, the coefficient of $s^nu^m$ in $\cC[G_C+s\,\dG;E_C+u\,\dE]$, is the sum over $\Gamma\in\cC$, over sets $\Lambda$ of $n$ scalar lines of $\Gamma$ and over sets $M$ of $m$ $E$-legs of $\Gamma$, of the value of $\Gamma$ with $\dG$ on $\Lambda$, $G_C$ on its other scalar lines, $\dE$ on $M$ and $E_C$ on its other $E$-legs.
\end{lemma}
\begin{proof}
A class value is a polynomial, multilinear in the scalar line kernels and in the $E$-legs (Wick expansion); put $s=u=1$. The counterterm coefficients are vertex data built from torus propagators at zero background \cite[(1.29)--(1.31)]{Ba3} and are not expanded: the replacement of propagators acts on lines only.
\end{proof}


\begin{definition}[cut skeletons and $\sigma$-classes]\label{def:a07:sigma}
Fix $\cC$ (type $T$, order $o$, labels $\cL$) and $(n,m)$. A \emph{cut skeleton} $\sigma$ consists of a partition $\cL=\cL_1+\dots+\cL_\rho$ into non-empty sub-multisets with orders $o_1+\dots+o_\rho=o$, $o_c\ge0$; a connected multigraph $Q$ on $\{1,\dots,\rho\}$ with $n$ edges (loops and multiple edges allowed); and an assignment of the external legs of $\cC$ to the components, each $E$-leg marked $\dE$ or $E_C$ with $m$ marked $\dE$. For each component $c$, $\cC_{\sigma,c}$ is the per-label full class of all connected configurations with labels $\cL_c$ and order $o_c$, whose external legs are the legs assigned to $c$ together with $\deg_Q(c)$ scalar \emph{stub} legs, with scalar lines $G_C$ and vector lines as in $\cC$, \emph{members without internal line and of order $0$ included}: equivalently, the coefficient of the monomial of the labels $\cL_c$ in the stub-Taylor expansion of the full connected generating function $\ln\int e^{-V(\cdot+\chi+\psi)}d\mu_{G_C}$, $\psi$ the stub field. The stub and $E$-legs may sit at any scalar slot of any vertex, differenced or not; the class sums over all of them. The \emph{$\sigma$-class} is the contraction $\Sigma_\sigma(\cC)\defeq\big\langle\bigotimes_c\cC_{\sigma,c},\bigotimes_{e\in Q}\dG_e\big\rangle$: an edge $\{c,c'\}$ contracts one stub of $c$ with one stub of $c'$ by $\dG$, a loop two stubs of $c$. It carries the Taylor weight $w_\sigma$, with $0\le w_\sigma\le C(\bar n)$.
\end{definition}

\begin{lemma}[regrouping into $\sigma$-classes]\label{lem:a07:R}
$\cC^{(n,m)}=\sum_\sigma w_\sigma E(\Sigma_\sigma(\cC))$, the sum over the cut skeletons of $(\cC;n,m)$.
\end{lemma}
\begin{proof}
\emph{The bijection.} To $(\Gamma,\Lambda,M)$ assign the components of $\Gamma\setminus\Lambda$ (connected through the kept lines, vector lines included, since $\Lambda$ consists of scalar lines), their label multisets and orders, the multigraph $Q$ that $\Lambda$ induces on them (connected, since $\Gamma$ is), the assignment of the external legs, and for each $c$ the restriction $\Gamma_c$ of $\Gamma$ to $c$ with the ends of $\Lambda$-lines at $c$ turned into stubs. Then $\Gamma_c\in\cC_{\sigma,c}$: it is connected, with labels $\cL_c$, order $o_c$ and its external legs; it may be a single vertex all of whose legs are stubs or external fields (when every line at that vertex is in $\Lambda$), or a path of order-$0$ vertices (when $\Lambda$ cuts out a stretch of an expanded line), which is why these members are in $\cC_{\sigma,c}$ (with \cite[p.~415]{Ba3}'s ``there is at least one internal line'' the map would not land in the product). Conversely, given $\sigma$, members $\Gamma_c\in\cC_{\sigma,c}$ and a pairing of all stubs compatible with $Q$, joining the paired stubs by lines gives a connected graph of type $T$, order $o$ and labels $\cL$, hence a member of the full class $\cC$, with $\Lambda$ the set of joining lines; the two maps are inverse. \emph{The weights.} In the labelled form with vertex labels $1,\dots,N$ and weight $\frac1{N!}\prod_v(-c_v)$, the components take disjoint sets of vertex labels, in $N!/\prod_cN_c!$ ways, so the weight is $\prod_c\frac1{N_c!}\prod_{v\in c}(-c_v)$ times the symmetry factor of identical components, parallel edges and loops of $Q$, which is that of the Taylor coefficients. \emph{Generating-function form} (the same identity, weights included): with $W(K;\chi)$ the connected generating function at fixed labels, scalar covariance $K$ and external field $\chi$, $W(G_C+s\,\dG;\chi)=\ln\E_{s\dG}\big[\exp W(G_C;\chi+\psi)\big]$, a formal Gaussian convolution in the stub field $\psi$ (the moments of covariance $K_1+K_2$ are the convolution of those of $K_1$ and $K_2$, an identity of polynomials in the kernels; no positivity of $\dG$ is used). By the linked-cluster theorem \cite{GJ} the right side is a sum of connected graphs whose vertices are the $\psi$-Taylor coefficients of $W(G_C;\chi+\psi)$, the classes $\cC_{\sigma,c}$ with their stubs, and whose lines are $\dG$; taking the coefficient of the labels and of $s^n$, with $\chi=E_C+u\,\dE$, and then of $u^m$, gives the claim. \emph{Products.} For a product of classes contracted by fixed many-variable fields \cite[p.~419]{Ba3}, apply the above factor by factor; the fixed fields stay external fields of the components they touch and join the components in step 5 below exactly as the $\dG$ fields do.
\end{proof}

The classes $\cC_{\sigma,c}$ are full per-label classes, not orbits: the finer groupings (\cref{lem:a07:orb}) fix slot data, which is needed where a bound depends on slot types; here only the $\norm\cdot_{1,\alpha_0}$ norms of $\dG$ and $\dE$ are consumed, which control both slot types (\cref{lem:a07:K3} has $\sigma,\sigma'\in\{0,D,HD\}$ in each variable), so the full grouping is admissible, and it is the form in which the renormalisation bound is available at a nonzero background. \Cref{lem:a07:R} is Ba{\l}aban's own step at \cite[p.~432]{Ba3} for the lines $\delta G_k$ (``we replace these lines by pairs of external legs \dots\ $G'$ may be unconnected \dots\ it can be represented as a sum of classes $G'_{\mathrm{ren}}$''), which for full classes is exact.

\begin{lemma}[the $\sigma$-classes at a nonzero background]\label{lem:a07:S}
For every cut skeleton $\sigma$ and component $c$:
\begin{enumerate}[label=\textup{(\alph*)},nosep,leftmargin=*]
\item $\cC_{\sigma,c}$ is a family of the form of \cite[Prop.~1]{Ba3} on the domain $\OmC$ at the background $\Bt$: vertices \cite[(1.6)--(1.15)]{Ba3} localised in $B^k(\Lsev)$, scalar lines $G_k(\OmC,\Bt)$, vector lines as printed, and external fields of several variables (the stub kernels $\dG$, ``treated as an external field'' \cite[pp.~414, 419]{Ba3}; the fields $\dE$, $E_C$; the remaining fields of $\cC$).
\item Ba{\l}aban's classes are sets of graphs built from the vertices of \cite[(1.5)]{Ba3}, and $\Bt$ enters only the line operators and the transporters of \cite[(1.12)--(1.15)]{Ba3}; so $G_{\mathrm{ren}}(G)$ is the same set of graphs at $\Bt=0$ and at $\Bt\ne0$. Every completion keeps type, order and label multiset; hence $G_{\mathrm{ren}}(G)\subset\cC_{\sigma,c}$ for every member $G$, and $\cC_{\sigma,c}$ is the union of the $G_{\mathrm{ren}}(G)$ over its members.
\item $\cC_{\sigma,c}$ is closed under the exchange of the two external scalar legs of every two-scalar-leg or $\phi\phi A$ subgraph, the whole graph included (\cref{lem:a07:full}(c)), exchanges that move a stub or change its slot type included; and its Taylor coefficients in the fluctuation of the background are again per-label full classes, every touched two-scalar-leg subgraph having its mirror.
\item Split $\cC_{\sigma,c}=\cC^{\mathrm{ren}}_{\sigma,c}+\cC^{0}_{\sigma,c}$, where $\cC^{\mathrm{ren}}_{\sigma,c}$ consists of the members of order $\ge1$ which are graphs in the sense of \cite[p.~415]{Ba3} together with the lone counterterm vertex \cite[(1.7)]{Ba3} (the counterterm member of the whole-graph pairing), and $\cC^0_{\sigma,c}$ of the bare single vertices with every leg a stub or an external field and of the order-$0$ paths of $n=0$ vertices joined by $G_C$ lines. Then $\cC^{\mathrm{ren}}_{\sigma,c}$ is a per-label full class of the form of \cite[Prop.~1]{Ba3}, to which the bound \cite[(1.33)]{Ba3} applies at block factor $5$ with constants ``independent of $\eps$, $k$, the domains $\Omega,\Omega_1,\Omega_2$, the vector field $\tilde B$'' \cite[p.~421]{Ba3} (\cref{thm:a06:BALA}; the unprinted steps of that proposition are those listed in Step~3 of the proof of \cref{thm:a06:T1}), with the localisation $h$ of each external field equal to $1$ on the label boxes of every vertex at which the field sits in some member of the class, the class being split by the cluster at which each leg sits (at most $V(\bar n)^{\#\mathrm{legs}}$ pieces). Each member of $\cC^0_{\sigma,c}$ obeys a bound of the form of \cite[(1.33)]{Ba3} directly: a bare vertex is at most its coupling times $\prod_i\norm{h_i\Phi_i}_{1,\alpha_0}$ summed over its cube or block (\cite[p.~421]{Ba3}: ``the vertex itself with all external legs \dots\ is treated separately''); an order-$0$ path is a tree whose only pieces of degree $\le0$ are the pairs $VG_CV$ of \cite[(2.4)]{Ba3}, so that \cite[Prop.~2.1]{Ba3} applies to it: each pair is treated by \cite[(2.8)--(2.9)]{Ba3} at one vertex, the derivative moved onto a stub, an $E$-leg or the external vector leg (allowed since the norm \cite[(1.32)]{Ba3} contains one derivative, and \cref{lem:a07:K3} covers the stub), after which all degrees are positive. No member of $\cC^0_{\sigma,c}$ is a partner that the bound needs: partners have the same order $\ge1$, and a zero-line member of order $\ge1$ other than a counterterm vertex is a (1.6) vertex with four external legs, alone in its class. The factors of the background in these bounds are $O(e(t)p)$ \cite[p.~412]{Ba3}.
\end{enumerate}
\end{lemma}
\begin{proof}
(a) By construction and \cref{lem:a07:C51}. (b) At \cite[p.~414]{Ba3} ``the vertices of (1.5) are given by (1.6)--(1.8), (1.10), (1.13), and (1.14) with $n'=0$ and $\tilde B=A^{(k)}$'', and the lines are ``$G_k(\Omega,\tilde B)$, $G_k(\Omega_2,\tilde B)$, $\delta G_k(\Omega,\Omega_2,\tilde B)$ or the operator (1.16)''; the completions (the counterterm graph of a divergent two-scalar-leg subgraph \cite[p.~420]{Ba3}, the permutation \cite[p.~429]{Ba3}, the two-vector-leg class in its maximal form \cite[p.~431]{Ba3}) keep type, order and labels (``localizations in these graphs should be in agreement''), and a per-label full class contains all graphs of its type, order and labels. (c) The first clause is \cref{lem:a07:full}(c); the second is the identity of generating functions in its proof: the expansion in the fluctuation $\tilde B'$ of the background is the Taylor expansion of the full connected generating function, whose coefficients are full classes with $\tilde B'$-legs carrying their partition labels, and the mirror of a touched two-scalar-leg subgraph is an exchange of the first clause. (d) The members of $\cC^{\mathrm{ren}}_{\sigma,c}$ form the per-label full class of \cite[Prop.~1]{Ba3}'s form by (a)--(c); the lone (1.7) vertex is the counterterm member of the whole-graph pairing (\cite[p.~417]{Ba3}: ``the same graph $G$ but with both external legs localized in $x$''). The localisation rule for a field that sits at several vertices of a class is the reading of \cite[p.~420]{Ba3} (``$h=1$ on $\mathrm{box}(v)$'') for classes; per piece the stub bound of step 4 below holds with the piece's $h$, whose supports are at least as far apart as the original ones. The bounds on $\cC^0_{\sigma,c}$ are as stated: a bare vertex carries no line and no loop; an order-$0$ path consists of $n=0$ vertices, each with at most two scalar legs, so it is a tree, and \cite[(2.8)]{Ba3} is the one-derivative move of \cite[Sect.~2]{Ba3} at a vertex of a (2.4)-type pair. The last two clauses are read off the orders.
\end{proof}

\begin{theorem}[\BALint: the interaction-term half of {\cite[Prop.~2.6]{Ba2}} at every background]\label{thm:a07:dom}
In the setting above, for every $\kappa>0$ there are $C_\kappa<\infty$ and $t_0(\kappa)>0$ such that for $t\le t_0(\kappa)$
\[
\sum_{\cL}\sum_{\cC}\sum_{(n,m)\ne(0,0)}\sum_{\sigma}w_\sigma\big|E\big(\Sigma^{(n,m)}_\sigma(\cC)\big)\big|\ \le\ C_\kappa\,t^{\kappa}\,|\Lsev|,
\]
the sums over every localisation $\cL$ of the stage (vertex cubes in $B^k(\Lsev)$, label boxes within $O(1)$ of it), every class $\cC$ with labels $\cL$, every $(n,m)\ne(0,0)$ and every cut skeleton $\sigma$. By \cref{lem:a07:M,lem:a07:R} the left side dominates $\sum_\cL\sum_\cC|\cC[G_A;E_A]-\cC[G_C;E_C]|$: this is \cite[Prop.~2.6]{Ba2} for the interaction terms, with error $O((L^k\eps)^\kappa)|\Lsev|$ for every $\kappa$, at every background $\Bt$, $\Bt=0$ included; $C_\kappa$ and $t_0(\kappa)$ depend on $\kappa$, on $(R,\varrho,b_0,\tilde p,\bar n,\alpha_0)$, on the couplings and on the constants of \textup{(P1)--(P3)} and of \cite[(1.33)]{Ba3} only. This proves \cref{hyp:k01:BALint}.
\end{theorem}

\begin{proof}
\emph{Step 1 (domain change only).} By \cref{lem:a07:C51} the replacement is $G_A\to G_C$ at the one background $\Bt$; $\OmC\subset B^k(\Lambda^{(k)}_0)\subset\OmA$, and $\Bt$ is regular on $\OmA$.

\emph{Step 2 (expansion and regrouping).} By \cref{lem:a07:M,lem:a07:R} the difference is $\sum_{(n,m)\ne(0,0)}\sum_\sigma w_\sigma E(\Sigma^{(n,m)}_\sigma(\cC))$.

\emph{Step 3 (depth).} Every vertex, every label box and every localisation support lies within $2+\frac14+\eta<3$ of $B^k(\Lsev)$, hence at depth $\ge D_7=5(r-1)-4$ in $\OmC$ by \cref{lem:a07:G}(a), with $D_7\ge R_0+2$ by \eqref{eq:a07:D7}. The domain hypothesis of \cite[Prop.~1]{Ba3} for $\Omega_2=\OmC$ (``the same conditions as $\Omega$'' \cite[p.~414]{Ba3}, including $\dist(\supp\tilde A,\partial\Omega_2)>2r$ \cite[p.~412]{Ba3}) holds for the $\tilde A$-legs actually present: with $\tilde A=\theta_{k+1}A'^{(k)}$ the raw support reaches $\partial\OmC$, but after \cite[p.~579]{Ba2} every vertex and every localised leg lies within $O(1)$ of $B^k(\Lsev)$, at depth $\ge D_7>2r$ (since $r\ge R\ge R_0+3>3$); multiplying $\tilde A$ by a cut-off equal to $1$ on that neighbourhood and supported at depth $>2r$ changes no term.

\emph{Step 4 (external fields).} For the localisations $h,h'$ of the two ends of an edge of $Q$ ($h=1$ on the label boxes of the vertices at which the stub can sit in its component, supported in their $\frac14$-neighbourhood), \cref{lem:a07:K3} with $(\Omega,\Omega_2,A)=(\OmA,\OmC,\Bt)$ gives
\[
\norm{h\,\dG\,h'}_{1,\alpha_0}\ \le\ C_hC_{\dG}\,e^{-2\delta_{\dG}D_7}\,e^{-\delta_{\dG}\dist(\supp h,\supp h')},
\]
with the mixed clauses $\sigma,\sigma'\in\{0,D,HD\}$ that the joining of components needs. For $\dE$: $\norm{h\,\dE}_{1,\alpha_0}\le2c_0e^{-\delta_0D_7}\norm f_\infty$, since (P1) (2.26) bounds $a_k\dG f_{\mathrm{in}}$ and (2.25) bounds $a_kG_Af_{\mathrm{out}}$ at points at distance $\ge D_7$ from $\supp f_{\mathrm{out}}\subset\OmC^c$. The norms of $E_C$ and of the other external fields are $\le Cp^a\lambda(t)^{-b}$ ((P1) (2.24)--(2.25) on $\OmC$ at depth $D_7$).

\emph{Step 5 (the renormalisation bound per $\sigma$-class).} By \cref{lem:a07:S}(d), \cite[(1.33)]{Ba3} or the direct bound applies to each $\cC_{\sigma,c}$ at fixed labels with the external fields of Step 4. The components are joined along a spanning tree of $Q$ by the decay of the stub kernels, the non-tree edges and loops entering as many-variable external fields (\cite[(3.2), p.~432]{Ba3}: ``the exponential factor in (1.33) is obtained from the estimates of external fields in (3.2)''). With $e_k,\lambda_k\le1$,
\[
\big|E\big(\Sigma^{(n,m)}_\sigma(\cC);\cL\big)\big|\ \le\ O(1)\,e^{-\delta''d(\cL)}\,\big(C'e^{-2\delta_{\dG}D_7}\big)^{n}\big(2c_0Cp\lambda^{-1/4}e^{-\delta_0D_7}\big)^{m}\big(Cp^a\lambda^{-b}\big)^{n_{\mathrm{ext}}},
\]
where $d(\cL)$ is the tree length of all localisation cubes over all components and $\delta''\defeq\frac12\min(\delta_0,\delta_{\dG},\delta_{\mathrm{ren}})$, $\delta_{\mathrm{ren}}$ the rate of \cite[(1.33)]{Ba3} (called $\delta_0$ there; it is not the $\delta_0$ of \textup{(P1)}): the trees inside the components come from (1.33), the edges between components from the stub decay, and the $\frac14$-neighbourhoods change the tree length by at most a constant times the number of cubes, a constant of $\bar n$ absorbed in $O(1)$.

\emph{Step 6 (size).} Since $n+m\ge1$ the product carries $e^{-\delta_*D_7}$ at least once, $\delta_*\defeq\min(2\delta_{\dG},\delta_0)$, and $D_7\ge r$; by \cref{lem:a07:G}(b), $p^a\lambda^{-b}e^{-\delta_*D_7}\le t^\kappa$ for $t\le t_0(\kappa)$, the background factors $O(e(t)p)$ of \cref{lem:a07:S}(d) being absorbed the same way.

\emph{Step 7 (sum).} Fix one cube in $B^k(\Lsev)$ ($|\Lsev|$ unit cubes, and every cube lies there); \cref{lem:a07:G}(c) sums the others with $e^{-\delta''d}$. The numbers of types, of orders $\le\bar n$, of pairs $(n,m)$ and of skeletons $\sigma$ are bounded by constants of $\bar n$ ($\rho\le V(\bar n)$, $n$ at most the number of scalar lines, $Q$ and the leg assignments finite).
\end{proof}

\begin{remark}\label{rem:a07:dom}
(i) The proof uses no telescoping in a fixed line order and no forest families: the terms with one difference line, grouped per line, carry their counterterm partners but not the two-vector-leg and permutation partners of the cut graph (\cref{rem:a07:orb}(iv)), and a lone one-loop two-vector-leg subgraph is $O(\eta^{-1})$ per graph, which $e^{-2\delta_{\dG}D_7}=O(t^\kappa)$ does not cure; the regrouping of \cref{lem:a07:R} restores the complete classes, at every background. (ii) The corresponding replacement in the polynomial $Q^{(k)}$ of \cite[p.~579]{Ba2}, the other half of \cite[Prop.~2.6]{Ba2}, is not used through \cref{hyp:k01:BALint}. (iii) On the torus skew in the rotated frame (the instance consumed by the rotation part through \cite[Thm.~3.1]{K1} there), \cref{lem:a07:M,lem:a07:R,lem:a07:S} are combinatorics without geometry, \cite[(1.33)]{Ba3} transfers as the other torus estimates do, and the proof of \cref{lem:a07:K3} uses no cubic symmetry, its inputs (P1)--(P3) being available on that torus: \cref{lem:a07:Ptorus,lem:a07:K3torus} below carry (P1)--(P3), \cref{lem:a07:K3} and \cref{thm:a07:dom} to every lattice torus fitted to the blocks, these tori included.
\end{remark}

% k0x_fill_skew: the inputs (P1)--(P3), the difference-kernel bound of lem:a07:K3 and \BALint\ on a lattice torus fitted to the blocks
% (input from sec/a07_k01steps_c.tex after rem:a07:dom). This is the transfer consumed by the rotation part on the tori of rot:hyp:F.

\subsection{The difference-kernel bound on a lattice torus fitted to the blocks}\label{ssec:a07:torus}

The rotation part applies King's \cite[Thm.~3.1]{K1}, and through it \BALint, on the tori $T_1=\DJ$ and $T_2=R^{-1}\DJ$ of \cref{rot:hyp:F}, each carrying both lattices $\Lk$ and $R^{-1}\Lk$ of \cref{rot:def:tori}(b). In the frame of the lattice $R^{-1}\Lk$ the torus $\DJ$ is the quotient of that lattice's own axis lattice by $R(50MJ\,\Z^3)$, and in the frame of $\Lk$ the torus $R^{-1}\DJ$ is the quotient by $R^{-1}(50MJ\,\Z^3)$: lattices of translations which are not rectangular in those frames. The statements of this appendix were written for subsets of $\eta\Z^3$. This subsection records that their proofs use the global geometry only through the two properties of \cref{def:a07:btorus}, which the tori of the rotation part have, so that \cref{lem:a07:CT,lem:a07:K3} and \cref{thm:a07:dom} hold on them with constants that do not depend on the torus.

\begin{definition}[a lattice torus fitted to the blocks]\label{def:a07:btorus}
In scale-$k$ units let $T=\eta\Z^3/\Gamma$, where $\Gamma\subset\Z^3$ is a lattice of rank $3$ (so $T$ carries the unit blocks and, since $\Gamma\subset\Z^3$, the blocks of every scale $j\le k$, of side $l_j=5^{j-k}$). Let $M_b\in\N$ be the side of a large block. $T$ is \emph{fitted to the blocks} if $\Gamma\subset M_b\Z^3$ and its sup-systole $\mathrm{sys}_\infty(T)\defeq\min\{|\gamma|_\infty:\gamma\in\Gamma\setminus\{0\}\}$ is at least $3M_b$. Distances on $T$ are the quotient sup-distances $\dist_T(x,x')\defeq\min_{\gamma\in\Gamma}|x-x'-\gamma|_\infty$; $\Omega\subset T$ denotes a union of large blocks, $T$ itself included, and ``regular on $\Omega$'' is \cite[(2.23)]{Ba1} read on the bonds of $\Omega$.
\end{definition}
If $T$ is fitted to the blocks in scale-$k$ units, so is its image in scale-$j$ units for every $j\le k$: the lattice becomes $l_j^{-1}\Gamma\subset l_j^{-1}M_b\Z^3\subset M_b\Z^3$ (as $l_j^{-1}=5^{k-j}\in\N$) and the systole grows by the factor $l_j^{-1}$. For the tori of the rotation part, in the frame of either lattice, the lattice of translations is $50MJ\,\Z^3$ or $R^{\pm1}(50MJ\,\Z^3)\subset10MJ\,\Z^3$ in physical units (\cref{rot:lem:geom}(b)). At Ba{\l}aban's step $k$ on the lattice of spacing $\eps_K=5^{-K}$ the scale-$k$ unit is $t=5^k\eps_K=5^{k-K}\le1$, so in scale-$k$ units these lattices lie in $10MJ\,5^{K-k}\,\Z^3$ and the systole is at least $10MJ\,5^{K-k}$. They are therefore fitted to the blocks at every step whenever the large-block side $M_b$ divides $10M$ and $3M_b\le10MJ$; in King's flow, the one run on these tori through \cite[Thm.~3.1]{K1}, the large block is King's and $M_b=M$.

\begin{lemma}[the printed inputs on a torus fitted to the blocks]\label{lem:a07:Ptorus}
Let $T$ be fitted to the blocks, $\Omega\subset T$ a union of large blocks and $A$ regular on $\Omega$. Then the bounds \textup{(P1)} (\cite[Prop.~2.1, (2.24)--(2.26)]{Ba1}, \cite[(1.9)--(1.12)]{Ba4}), \textup{(P2)} (\cite[Props.~2.2--2.3]{Ba1}, \cite[(1.13)--(1.20)]{Ba4}) and the positivity \cite[(1.8)]{Ba4} hold for the operators $G_k(\Omega,A)$, $\Delta^{(k)}(\Omega,A)$, $C^{(k)}_\Lambda(\Omega,A)$ on $T$, with the distances of \cref{def:a07:btorus}, with constants $c_0$, $\delta_0$, $R_0$, $c_2$, $\theta_0$, $\gamma$ that depend on $d$, $a$, $M_b$ and $L$ ($c_0$ also on $\alpha$) and not on $\Gamma$, $k$, $\Omega$, $\Lambda$ or $A$. For $\Omega=T$ the restriction to points at distance $\ge R_0$ from $\Omega^c$ is void. The same holds in scale-$j$ units for every $j\le k$.
\end{lemma}
\begin{proof}
Ba{\l}aban's proofs \cite[Sect.~2, pp.~574--586]{Ba4} are written for a union $\Omega$ of large blocks in $\eta\Z^d$ (``the torus identified with a rectangular parallelepiped''). They use the geometry of the ambient lattice in exactly the following places, each of which we verify on $T$.

\emph{(1) The covering cubes.} Let $v_j$ run through the vertices of the large-block grid of $T$ and $C_j\defeq v_j+[-M_b,M_b)^3$, the union of the $2^3$ large blocks with vertex $v_j$ (\cite[p.~578]{Ba4}: ``a cube with center $Mj$ and size equal to $2M$''). Since $2M_b<\mathrm{sys}_\infty(T)$, the projection $\eta\Z^3\to T$ is injective on each lift of $C_j$; it preserves bonds, unit blocks and the contours $\Gamma^{(k)}_{y,x}$ of the averaging operators, which lie inside blocks \cite[(1.4)]{Ba4}. Hence $\Pi_j\defeq\Omega\cap C_j$ is a union of at most $2^3$ large blocks in an embedded cube, the constant gauge $\lambda_0(x)=A_0\cdot(x-c_j)$ is defined on $C_j$, and \cite[Lemmas~2.1, 2.2, 2.4]{Ba4}, which are statements about $G_k(\Pi_j,A_j)$ on a parallelepiped of $\eta\Z^d$ with a local field $A_j$ built from $A|_{C_j}$, apply with their constants, which depend on $d$, $M_b$, $\alpha$ and the bare $a$ only.

\emph{(2) The partition of unity.} Let $h\in C^\infty_c((-\frac23,\frac23))$ with $h=1$ on $[-\frac13,\frac13]$ and $\sum_{n\in\Z}h(\cdot-n)^2=1$, and $h_j(x)\defeq\prod_\mu h((x_\mu-(v_j)_\mu)/M_b)$ \cite[p.~575]{Ba4}. Because $\Gamma\subset M_b\Z^3$, the family $\{h_j\}$ is $\Gamma$-periodic, descends to $T$, and $\sum_jh_j^2=1$ on $T$ with $\supp h_j\subset C_j$; so \cite[(2.2)--(2.9)]{Ba4} hold on $T$ as on $\eta\Z^d$ (if $\Omega=T$, then $\partial\Omega=\emptyset$ and (2.6) is trivial).

\emph{(3) The counts.} Each $j$ has at most $3^3$ neighbours and each point lies in at most $2^3$ of the supports, so the remainder of the parametrix has $\norm R_{2,2}\le c_dO(M_b^{-1})<1$ as on \cite[p.~577]{Ba4}, and the number $2^d(3^d)^{n-1}2^d$ of walks of length $n$ \cite[p.~579]{Ba4} is unchanged.

\emph{(4) The metric.} Consecutive cubes of a walk have centres at quotient distance $\le M_b$, so a walk of length $n$ from a cube meeting $\supp f$ to a cube containing $x$ has $n\ge\dist_T(x,\supp f)/M_b-2$; the geometric series \cite[(2.22)]{Ba4} then gives $\delta_0$, $c_0$ as on $\eta\Z^d$. The definition \cite[(2.19)]{Ba4} of $R_0$ and the implication ``$\dist_T(\{x,x'\},\Omega^c)\ge R_0$ implies that the first $n_0$ cubes of the walk lie in $\Omega$'' use only this. The cancellation of the interior-cube terms in $G_k(\Omega,A)-G_k(\Omega_0,A)$, which gives the two-sided factor of \cite[(2.26)]{Ba1}, uses (1) and (2). If $\Omega=T$, every $\Pi_j=C_j$ is a full cube and no restriction on $x$, $x'$ is needed \cite[p.~611]{Ba1}.

\emph{(5) Positivity.} For \cite[(1.8)]{Ba4}, $P_k(A)$ is block-diagonal, so dropping the bonds between unit blocks bounds the form from below by a sum over the unit blocks $\Delta$ of $\langle\phi,(-\Delta^N_{A,\Delta}+a_kP_k(A))\phi\rangle_\Delta$; on each block the gauge of (1) and the expansion in the field reduce the bound to $A=0$ \cite[p.~580]{Ba4}, where it is the Poincar\'e inequality on a unit block, which embeds in $T$. So $\gamma$ is Ba{\l}aban's.

The block operators of (P2) are built from $G_k(\Omega,A)$ and the averaging operators, and the proofs of \cite[Props.~2.2--2.3]{Ba1}, \cite[Sect.~3]{Ba4} use the same cubes, partition and walk and no further property of the ambient lattice; the restricted covariances $C^{(k)}_\Lambda$ for a union $\Lambda$ of large blocks are handled by the same walk on $\Lambda$. Nothing else in \cite[Sects.~2--3]{Ba4} refers to the global geometry. The last sentence of the lemma follows because the rescaled torus is fitted to the blocks (the remark after \cref{def:a07:btorus}) and the statements are the same at every scale.
\end{proof}

\begin{lemma}[\cref{lem:a07:CT,lem:a07:K3} and \BALint\ on a torus fitted to the blocks]\label{lem:a07:K3torus}
Let $T$ be fitted to the blocks. Then:
\begin{enumerate}[label=\textup{(\alph*)},nosep,leftmargin=*]
\item \cref{lem:a07:CT} holds for subsets of $T$ with the quotient metric, with the same constants.
\item \cref{lem:a07:K3} holds for unions $\Omega_2\subset\Omega$ of large blocks of $T$ and $A$ regular on $\Omega$, with $\dist_T$ in place of $\dist$ and with $C_{\dG}$, $\delta_{\dG}$ depending on $d$, $a$, $M_b$, $L$ and $\alpha$ only.
\item \cref{thm:a07:dom} holds on $T$, with the propagator domains $\OmA\supset\OmC$ and the background $\Bt$ of Ba{\l}aban's stage on $T$ (\cite[Sect.~2]{Ba2} run on $T$, as in \cref{thm:a06:BALA}) and with the renormalisation bound \cite[Prop.~1, (1.33)]{Ba3} on $T$; its constants $C_\kappa$, $t_0(\kappa)$ are those of the statement on $\eta\Z^3$ and do not depend on $\Gamma$.
\end{enumerate}
In particular all three hold on the tori $T_1$, $T_2$ of \cref{rot:hyp:F}, in the frame of either lattice.
\end{lemma}
\begin{proof}
(a) The proof of \cref{lem:a07:CT} uses the weight $\rho=\dist(\cdot,\supp f')$ only through two properties: $\rho$ is $1$-Lipschitz along bonds (the ends of a bond are at distance $\eta$), and each site is the final point of $d$ bonds. Both hold on $T$ with $\rho\defeq\dist_T(\cdot,\supp f')$, since $\dist_T(\cdot,S)$ is $1$-Lipschitz for the quotient metric and the bond structure is that of $\eta\Z^3$ locally. The lower bound $T\ge\gamma$ it needs is \cite[(1.8)]{Ba4} on $T$ (\cref{lem:a07:Ptorus}). The constants are unchanged.

(b) The proof of \cref{lem:a07:K3} consists of the slice decomposition $G_k(\Omega,A)=G^{[k]}_1+\sum_{j=1}^{k-1}S_j$, an operator identity which needs only that the blocks of every scale $j\le k$ partition $\Omega$ (true on $T$: a union of large blocks is a union of unit blocks of every scale $j\le k$, and these descend to $T$ because $\Gamma\subset M_b\Z^3\subset l_j\Z^3$); the rescaling to scale-$j$ units, under which $T$ stays fitted to the blocks; the slice bounds, which are (P1) at scale $j$ for the legs, (P2) at scale $j$ for the block covariances and (P3) at scale $1$, all available on $T$ with $\Gamma$-free constants by \cref{lem:a07:Ptorus} and (a); the sums over the block centres $\hat y,\hat y'$, which use \cref{lem:a07:G}(c) in the form $\sum_{\square}e^{-\delta\dist_T(\square,S)}\le C(\delta)|S|$, true on $T$ because the projection of the cubes of $\eta\Z^3$ onto those of $T$ is onto, so the sum on $T$ is at most the sum on $\eta\Z^3$ over a lift of $S$; and the triangle inequality for the metric. Every step is therefore available on $T$ with the quotient metric, and the proof goes through word for word. The constants $C_{\dG}$, $\delta_{\dG}$ are built from those of (P1)--(P3), $L$ and $\alpha$.

(c) In the proof of \cref{thm:a07:dom}: \cref{lem:a07:M,lem:a07:R,lem:a07:S} are statements about the expansion and regrouping of graph sums and about the sets of graphs forming a class; they do not refer to the geometry. \cref{lem:a07:C51} uses only that the three operators depend on the field through the bonds of their domain and on contours inside blocks. \cref{lem:a07:G}(a) uses the construction of the layers $\Lambda^{(k)}_i$ by sup-neighbourhoods of unions of large blocks and sup-metric geodesics, which exist in the quotient metric (a shortest path lifts to $\eta\Z^3$); (b) is a scalar inequality; (c) holds on $T$ as shown in (b) above, the number of cubes of $T$ at a given tree distance from a fixed cube being at most the number on $\eta\Z^3$. Steps 1--3 of the proof use (a) of \cref{lem:a07:G} and the definitions of the stage; Step 4 uses (b) above and (P1) on $T$; Step 5 uses \cite[Prop.~1, (1.33)]{Ba3} on $T$, which is Ba{\l}aban's renormalisation bound for a torus aligned with the blocks (\cref{thm:a06:BALA}); Steps 6 and 7 use \cref{lem:a07:G}(b), (c). The constants are those of the statement on $\eta\Z^3$.

For the tori of the rotation part, both lattices of translations are contained in $10MJ\,\Z^3$ in the frame in question (\cref{rot:lem:geom}(b)), so the tori are fitted to the blocks with $M_b=M$ by the remark after \cref{def:a07:btorus}.
\end{proof}

\begin{remark}\label{rem:a07:torus}
(i) The two conditions of \cref{def:a07:btorus} are what Ba{\l}aban uses of the rectangular geometry: that the large blocks tile the torus at every scale (the partition of unity and the block operators), and that a cube of two large blocks' side embeds (the local parametrix and the local gauge). A rectangular torus of sides that are multiples of $M_b$ and at least $3M_b$ satisfies both. (ii) The systole condition can be removed by passing to the covering torus $\eta\Z^3/3\Gamma$ and pulling back, at the price of the factor $3^3$ in the kernel bounds from the image sum over the deck group; it is not needed for the tori of the rotation part, whose systole is $10MJ$ at least.
\end{remark}

